% from ega1-9-fr.tex
\oldpage[I]{173}

\item[(ii)] De nouveau $X$ est réunion finie d'ouverts affines $U_i$
($1\leq i\leq r$) tels que $\mathscr{L}|U_i\simeq\mathscr{O}_X|U_i$, et pour
chaque $i$,
$(s\otimes f^{\otimes n})|(U_i\cap X_f)$ s'identifie par l'isomorphisme
précédent à
$(f|(U_i\cap X_f))^n(s|(U_i\cap X_f))$. On sait alors
(\hyperref[I.1.4.1-fr]{1.4.1}) qu'il existe un entier $n>0$ tel que pour
chaque $i$, $(s\otimes f^{\otimes n})|(U_i\cap X_f)$ se prolonge en une
section $s_i$ de $\mathscr{F}\otimes\mathscr{L}^{\otimes n}$ au-dessus de
$U_i$. Soit $s_{i|j}$ la restriction de $s_i$ à $U_i\cap U_j$~; on a par
définition $s_{i|j}-s_{j|i}=0$ dans $X_f\cap U_i\cap U_j$. Or, si $X$ est un
espace noethérien, $U_i\cap U_j$ est quasi-compact~; si $X$ est un schéma,
$U_i\cap U_j$ est un ouvert affine
(\hyperref[I.5.5.6-fr]{5.5.6}), donc encore quasi-compact. En vertu de (i),
il existe donc un entier $m$ (indépendant de $i$ et $j$) tel que
$(s_{i|j}-s_{j|i})\otimes f^{\otimes m}=0$. On en conclut aussitôt qu'il
existe une section $s'$ de
$\mathscr{F}\otimes\mathscr{L}^{\otimes(n+m)}$ au-dessus de $X$, induisant
$s_i\otimes f^{\otimes m}$ au-dessus de chaque $U_i$, et induisant par suite
$s\otimes f^{\otimes(n+m)}$ au-dessus de $X_f$.
\end{enumerate}
\end{proof}

Les corollaires qui suivent donnent une interprétation du th. (9.3.1) en un
langage plus algébrique~:

\begin{corollary}[9.3.2]
\label{I.9.3.2-fr}
Les hypothèses étant celles de (9.3.1), considérons l'anneau gradué
$A_* = \Gamma_*(\mathscr{L})$ et le $A_*$-module gradué
$M_* = \Gamma_*(\mathscr{L},\mathscr{F})$
(\hyperref[0.5.4.6-fr]{0, 5.4.6}). Si $f\in A_n$, où $n\in\mathbb{Z}$, alors
on a un isomorphisme canonique
\[
  \Gamma(X_f,\mathscr{F})\xrightarrow{\sim}((M_*)_f)_0
\]
(sous-groupe du module de fractions $(M_*)_f$ formé des éléments de degré
$0$).
\end{corollary}

\begin{corollary}[9.3.3]
\label{I.9.3.3-fr}
On suppose vérifiées les hypothèses de (9.3.1), et on suppose en outre que
$\mathscr{L}=\mathscr{O}_X$. Alors si on pose
$A=\Gamma(X,\mathscr{O}_X)$, $M=\Gamma(X,\mathscr{F})$, le $A_f$-module
$\Gamma(X_f,\mathscr{F})$ est canoniquement isomorphe à $M_f$.
\end{corollary}

\begin{proposition}[9.3.4]
\label{I.9.3.4-fr}
Soient $X$ un préschéma noethérien, $\mathscr{F}$ un
$\mathscr{O}_X$-Module cohérent, $\mathscr{J}$ un faisceau cohérent d'idéaux
dans $\mathscr{O}_X$ tels que le support de $\mathscr{F}$ soit contenu dans
celui de $\mathscr{O}_X/\mathscr{J}$. Alors il existe un entier $n>0$ tel que
$\mathscr{J}^n\mathscr{F}=0$.
\end{proposition}

\begin{proof}
Comme $X$ est réunion finie d'ouverts affines dont les anneaux sont
noethériens, on peut supposer $X$ affine d'anneau $A$ noethérien~; alors
$\mathscr{F}=\widetilde M$, où $M=\Gamma(X,\mathscr{F})$ est un $A$-module
de type fini, et $\mathscr{J}=\widetilde{\mathfrak{J}}$, où
$\mathfrak{J}=\Gamma(X,\mathscr{J})$ est un idéal de $A$
(\hyperref[I.1.4.1-fr]{1.4.1} et \hyperref[I.1.5.1-fr]{1.5.1}). Comme $A$
est noethérien, $\mathfrak{J}$ admet un système fini de générateurs $f_i$
($1\leq i\leq m$). Par hypothèse, toute section de $\mathscr{F}$ au-dessus
de $X$ est nulle dans chacun des $D(f_i)$~; si $s_j$ ($1\leq j\leq q$) sont
des sections de $\mathscr{F}$ engendrant $M$, il existe donc un entier $h$
indépendant de $i$ et $j$, tel que $f_i^h s_j=0$
(\hyperref[I.1.4.1-fr]{1.4.1}), donc $f_i^h s=0$ pour tout $s\in M$. On en
conclut que si $n=mh$, on a $\mathfrak{J}^nM=0$, et par suite le
$\mathscr{O}_X$-Module correspondant
$\mathscr{J}^n\mathscr{F}=\widetilde{\mathfrak{J}^nM}$
(\hyperref[I.1.3.13-fr]{1.3.13}) est nul.
\end{proof}

\begin{corollary}[9.3.5]
\label{I.9.3.5-fr}
Sous les hypothèses de (9.3.4), il existe un sous-préschéma fermé $Y$ de
$X$, dont l'espace sous-jacent est le support de
$\mathscr{O}_X/\mathscr{J}$, et tel que, si $j:Y\to X$ est l'injection
canonique, on ait $\mathscr{F}=j_*(j^*(\mathscr{F}))$.
\end{corollary}

\begin{proof}
Notons d'abord que les supports de $\mathscr{O}_X/\mathscr{J}$ et de
$\mathscr{O}_X/\mathscr{J}^n$ sont les mêmes car si
$\mathscr{J}_x=\mathscr{O}_x$, on a aussi
$\mathscr{J}_x^n=\mathscr{O}_x$, et on a d'autre part
$\mathscr{J}_x^n\subset\mathscr{J}_x$ pour tout $x\in X$. On peut donc en
vertu de (9.3.4) supposer que $\mathscr{J}\mathscr{F}=0$~; on prend alors
pour $Y$ le sous-préschéma fermé de $X$ défini par $\mathscr{J}$, et comme
$\mathscr{F}$ est alors un $(\mathscr{O}_X/\mathscr{J})$-Module, la
conclusion est immédiate.
\end{proof}

